% ============================================================
% GROUND — Fourth Layer: The Instrument Built from This
% ============================================================
% Technical description of the work as research instrument:
%   — Substrate boxes: materials, construction, basin design
%   — Piezo wiring, sensitivity calibration thresholds
%   — Electret microphone placement
%   — Max/MSP patch: signal routing, filtering, granular chain
%   — Spatialization system: multi-channel, migration-data-informed
% Agential Realism (Barad): the substrate box as apparatus.
% The apparatus produces the phenomenon through intra-action —
% not just measuring what is already there.
% ST Cycle applied fully to this instrument:
%   Unsettling → Remixing → Embodying → Transmitting
% ST Domains applied:
%   FS (speculative recognition) / MM (earth as sedimented memory)
%   SW (multisensory, touch-to-sound interface)
%   IW (can participants sense the historical weight?)
% What remains unresolved; productive failures; next ST cycle.
%
% Key texts:
%   Barad — Meeting the Universe Halfway: apparatus,
%   intra-action, the phenomenon produced through entanglement.
%   Guerra & Ostergaard — Technopoiesis: the emergent
%   feedback between cultural logic and material form.
% ============================================================

\section*{Fourth Layer: The Instrument Built from This}

\ifdraft
\lipsum[1-4]
\fi
